% SPDX-FileCopyrightText: 2026 Will Cook
% SPDX-License-Identifier: CC-BY-4.0
%
% One of the Erdős Problem Notes: a short standalone treatment of a single
% problem in the ErdosProblems expansion library.  Build with
% `tectonic erdos-243-reciprocal-tail-rigidity.tex` or `make`.
\documentclass[11pt]{article}
% SPDX-FileCopyrightText: 2026 Will Cook
% SPDX-License-Identifier: CC-BY-4.0
%
% Shared preamble for the Erdős Problem Notes: the short, standalone,
% problem-owned manuscripts covering the expansion library `ErdosProblems`.
%
% One note owns one Erdős problem.  Each note repeats whatever context its own
% reader needs, so that a reader who arrives at a single problem never has to
% assemble it from the gateway paper.  Deliberate repetition across notes is the
% design, not an oversight.
%
% Source links resolve against one pinned formal-source commit, declared once
% below.  `scripts/check_problem_note_sources.py` verifies that every authored
% (file, line, declaration) triple names that exact declaration in the pinned
% snapshot, so a link cannot silently rot as later work moves lines.

\usepackage{paper-house-style}
\usepackage{array}
\usepackage{booktabs}
\usepackage{enumitem}
\setlist{topsep=0.35em,itemsep=0.2em}
\usepackage[colorlinks=true,linkcolor=linkinternal,urlcolor=linkexternal,
  citecolor=linkinternal]{hyperref}
\hypersetup{bookmarksnumbered=true,bookmarksopen=true,bookmarksopenlevel=1,
  pdfauthor={Will Cook},pdflang={en-GB}}

% ---------------------------------------------------------------------------
% Pinned formal source.  \commit names the checkpoint every link resolves
% against; the checker reads that snapshot rather than the working tree, so the
% printed coordinates stay correct however later waves move declarations.
% ---------------------------------------------------------------------------
\newcommand{\commit}{e5fd26a6d1b1a346430205eab5385b4c9565d004}
\newcommand{\commitshort}{e5fd26a6d1b1}
\newcommand{\repobase}{https://github.com/wcook04/plectis-lean-erdos249-257/blob/\commit}
\newcommand{\sourceurl}{https://github.com/wcook04/plectis-lean-erdos249-257/tree/\commit}
\newcommand{\PX}{ErdosProblems}

% Declaration names stay inside link targets.  The printed label is either a
% spaced, lower-case reading of the declaration name (\lref) or a short authored
% phrase (\lword); neither prints a module path or a raw Lean identifier.
\ExplSyntaxOn
\cs_new_protected:Npn \noteleanlabel #1
  {
    \tl_set:Nx \l_tmpa_tl { \tl_to_str:n {#1} }
    \regex_replace_all:nnN { _+ } { \c{space} } \l_tmpa_tl
    \regex_replace_all:nnN
      { ([a-z0-9])([A-Z]) }
      { \1\c{space}\2 }
      \l_tmpa_tl
    \tl_set:Nx \l_tmpa_tl { \tl_lower_case:n { \tl_use:N \l_tmpa_tl } }
    \tl_use:N \l_tmpa_tl
  }
\ExplSyntaxOff
\newcommand{\leanlabel}[1]{\noteleanlabel{#1}}
\newcommand{\lref}[3]{\href{\repobase/\PX/#1\#L#2}{\leanlabel{#3}}}
\newcommand{\lword}[4]{\href{\repobase/\PX/#1\#L#2}{#4}}
\newcommand{\lloc}[2]{\href{\repobase/\PX/#1\#L#2}{Lean source}}

% Links into a sibling library, named repository-relative.  The expansion
% library is implicit in \lref/\lword, because that is where problem-owned
% work lands.  Several problems have their headline theorem in the older
% reviewed #249/#257 corpus, and a note that could not name that library
% would have to paraphrase a checked statement instead of linking it.
\newcommand{\mref}[3]{\href{\repobase/#1\#L#2}{\leanlabel{#3}}}
\newcommand{\mword}[4]{\href{\repobase/#1\#L#2}{#4}}
\newcommand{\mloc}[2]{\href{\repobase/#1\#L#2}{Lean source}}
\emergencystretch=2em

\newtheorem{theorem}{Theorem}[section]
\newtheorem{proposition}[theorem]{Proposition}
\newtheorem{lemma}[theorem]{Lemma}
\newtheorem{corollary}[theorem]{Corollary}
\theoremstyle{definition}
\newtheorem{definition}[theorem]{Definition}
\newtheorem{example}[theorem]{Example}
\newtheorem{problem}[theorem]{Problem}
\theoremstyle{remark}
\newtheorem*{remark}{Remark}

\newcommand{\N}{\mathbb{N}}
\newcommand{\Npos}{\mathbb{N}_{>0}}
\newcommand{\Z}{\mathbb{Z}}
\newcommand{\Q}{\mathbb{Q}}
\newcommand{\R}{\mathbb{R}}
\newcommand{\lcm}{\operatorname{lcm}}
\newcommand{\ph}{\varphi}

% The series line printed under every note's title block.
\newcommand{\seriesline}{Erd\H{o}s Problem Notes}

% Production provenance.  The wording is fixed by the corpus provenance
% contract and reproduced verbatim in every manuscript in this repository, so
% the papers cannot drift into different accounts of how they were made.
\newcommand{\noteprovenance}{\provenancenote{\emph{Authorship and AI use.} The
prose, formal proofs, and software in this version were drafted and revised by
large-language-model agents under Will Cook's direction.  Cook set the
objectives and acceptance criteria, reviewed and authorised the manuscript, and
is responsible for its contents.  The agents are production tools, not authors.
See \emph{Statements and declarations}.}}

% The status paragraph every note carries, in its own words, before any result.
% Kernel checking establishes the exact Lean proposition; it does not establish
% that the proposition is the interesting one, that it is new, or that the
% Erdős problem is closed.
\newcommand{\statusboundary}{%
  \emph{Status.}  The problem treated here is open, and this note does not
  close it.  Every statement below marked as checked is a proposition that the
  pinned Lean kernel accepts from the sources this note links to, with no
  \texttt{sorry}, no added axiom, and no unchecked evaluation.  That is a claim
  about the formal statement, not about its mathematical interest, its
  novelty, or the original problem.  The unresolved obligations are named
  exactly, in their own section, and none of the finite computations,
  reductions, or no-go results here removes one of them.}

% Problem-specific source pin: #243's cancellation and feedback interfaces
% advance independently of the corpus default.
\renewcommand{\commit}{3d6d938d696fed0fb71dd55115a18a73738ff223}
\renewcommand{\commitshort}{3d6d938d696f}
\hypersetup{
  pdftitle={Excluding the bounded negative part: a bounded-defect theorem for Erdős Problem 243},
  pdfsubject={Erdős Problem Notes: Problem 243},
  pdfkeywords={irrationality, Ahmes series, Sylvester sequence, unit fractions, Lean 4},
}

% Names for the state maps and the two search objects.  They are set as
% upright operators so that they read as maps rather than as products.
\newcommand{\sylnext}{\operatorname{syl}}
\newcommand{\den}{\operatorname{den}}
\newcommand{\tail}{\operatorname{tail}}
\newcommand{\ctr}{\operatorname{ctr}}
\newcommand{\defect}{\operatorname{def}}
\newcommand{\fnum}{\operatorname{num}}
\newcommand{\surv}{\operatorname{Surv}}

\runninghead{Erd\H{o}s \#243: excluding the bounded negative part}
\title{Excluding the Bounded Negative Part}
\subtitle{A bounded-defect theorem for Erd\H{o}s Problem~\#243}
\author{Will Cook}
\date{July 2026; second review revision, 16 September 2026}

\begin{document}
\maketitle
\noteprovenance

\begin{abstract}
Let $a_1<a_2<\cdots$ be positive integers with $a_{n+1}/a_n^2\to1$
and rational reciprocal sum, and put $P_n=\prod_{j<n}a_j$. We prove that
\[
 \limsup_{n\to\infty}\frac{P_n}{a_n}
 \left(\frac{a_n^2}{a_{n+1}}-1\right)<+\infty
\]
forces $a_{n+1}=a_n^2-a_n+1$ eventually. We use Koizumi's integral
rational-tail recurrence, extending eventual nonnegative error to an
arbitrarily signed error with bounded negative part. The proof combines
integer descent, stabilisation of a common divisor and a first crossing of
a Chinese-remainder block. A second argument gives an LCM-weighted record
criterion. These are sufficient criteria: neither the required defect bound
nor the record budget is deduced from the unrestricted hypotheses.
\end{abstract}

\section{Introduction}\label{sec:problem}
\paragraph{Conventions and reading route.}
Write $z_+=\max(z,0)$, take empty products and least common multiples to be
$1$, and allow deletion of a finite prefix. In a higher-order rate such as
$1+3/n+o(n^{-3})$, retain the original index $n$: deleting terms does not
authorise renumbering that rate.
A \emph{Sylvester tail} means eventual equality
$a_{n+1}=a_n^2-a_n+1$, not a finite reciprocal expansion.
For the main argument read the next corollary, the integer-state theorem,
Section~\ref{sec:bounded} and the transfer in Section~\ref{sec:transfer}.
Sections~\ref{sec:mass}--\ref{sec:lcmrecords} give alternative criteria;
Section~\ref{sec:secondaryrate} records ordinary consequences whose full
proofs are in the companion reasoning paper.

For $a_n>1$, the recurrence $a_{n+1}=a_n^2-a_n+1$ is singled out by the telescope
\[
 \frac1{a_n-1}=\frac1{a_n}+\frac1{a_{n+1}-1}.
\]
Thus a Sylvester tail has the exact value $1/(a_n-1)$ beginning at index
$n$.  The question is whether rationality and the asymptotic relation
$a_{n+1}\sim a_n^2$ force that exact recurrence.  The following result
isolates a sufficient one-sided bound on the departure from it.

\begin{corollary}[original-coordinate bounded defect]\label{res:originalbounded}
Let $a_1<a_2<\cdots$ be positive integers, $a_{n+1}/a_n^2\to1$, and
$\sum_{n\ge1}1/a_n\in\Q$.  Put $P_n=\prod_{j<n}a_j$.  If
\[
 \limsup_{n\to\infty}\frac{P_n}{a_n}
 \left(\frac{a_n^2}{a_{n+1}}-1\right)<+\infty,
\]
then $a_{n+1}=a_n^2-a_n+1$ for all sufficiently large $n$.
\end{corollary}

The upper limit is allowed to be $-\infty$; the condition means an eventual
finite upper bound, not a bound on the absolute value. No monotonicity or
summability of the growth defect is assumed.

\noindent The proof is given in Section~\ref{sec:transfer}. The full
original-coordinate implication, including the canonical analytic transfer,
is also formalised as \href{https://github.com/wcook04/plectis-erdos/blob/3d6d938d696fed0fb71dd55115a18a73738ff223/lean/ErdosProblems/Erdos243/PaperCompleteR7/ProductDefect.lean\#L211}{the original-coordinate bounded-defect theorem}. Its hypothesis is an eventual upper bound on the displayed
product defect, equivalent here to the finite-upper-limit condition. The
integer-state theorem below isolates the arithmetic step; it is not the
whole extent of the formalisation.

\paragraph{The obstruction.}
Clear the rational reciprocal tail as $C_n/D_n$; these are the
denominator-cleared tail variables of Koizumi's pseudo-greedy formulation
\cite[Corollary~3, p.~9; Lemma~4, pp.~11--12]{koizumi2025}.  On a Sylvester tail,
$D_n=(a_n-1)C_n$; the integer error $E_n=D_n-(a_n-1)C_n$ measures the
failure of this identity.  Its update $C_{n+1}=C_n-E_n$ turns a bound on
negative error into a bound on upward increments.  Small increments alone
are insufficient.  The arithmetic input is that, after the common divisor
stabilises, every old multiplier must remain coprime to every later
reduced numerator.  The Chinese remainder theorem then produces a block of
forbidden heights that no permitted jump can cross.  A numerator tending
to infinity would have to cross it.  Section~\ref{sec:lcmrecords} adapts this crossing
argument when only jumps setting new records are charged.

\begin{samepage}
Write $\ctr(a,D,C)=D-(a-1)C$. The following integer-state theorem
implies Corollary~\ref{res:originalbounded}; the ordinary tail transfer is
proved in Section~\ref{sec:transfer}. Its centring hypothesis is redundant
under normalised vanishing and is retained to match the source statement.

\begin{theorem}[bounded negative part]\label{res:bounded}
Let $a,C,D:\N\to\N$ and $E:\N\to\Z$ satisfy
\begin{enumerate}
\item $a_n>1$ and $C_n>0$ for every $n$;
\item the exact dynamics $C_{n+1}+D_n=a_nC_n$ and $D_{n+1}=a_nD_n$;
\item $E_n=\ctr(a_n,D_n,C_n)$ for every $n$;
\item \emph{eventual strict centring}: $|E_n|<C_n$ for all large $n$;
\item \emph{eventually bounded negative part}: $-B\le E_n$ for all large $n$,
for some integer $B\ge0$;
\item \emph{normalised vanishing}: for every integer $K\ge1$ there is an $N$ with
$K\,|E_n|<C_n$ for all $n\ge N$.
\end{enumerate}
Then $E_n=0$ for all sufficiently large $n$: this is the
\lword{Erdos243/ReciprocalTailRigidity.lean}{2360}{eventuallyBoundedNegativePart_eventually_zero}{eventual bounded-negative rigidity}.
\end{theorem}
The Lean proof shifts the orbit past the centring and lower-bound thresholds
and applies the
\lword{Erdos243/ReciprocalTailRigidity.lean}{2271}{boundedNegativePart_eventually_zero}{bounded-negative-part rigidity},
with centring and bound at every index.
\end{samepage}

\begin{problem}[Erd\H{o}s \#243]\label{res:problem}
Let $1\le a_1<a_2<\cdots$ be a sequence of integers with
\[
 \lim_{n\to\infty}\frac{a_n}{a_{n-1}^{2}}=1
 \qquad\text{and}\qquad
 \sum\frac{1}{a_n}\in\Q .
\]
Then $a_n=a_{n-1}^{2}-a_{n-1}+1$ for all sufficiently large $n$.
\end{problem}

The question is recorded by Erd\H{o}s and Graham~\cite{erdosgraham1980} and by
Erd\H{o}s~\cite[p.~105]{erdos1988}, and Bloom's catalogue lists it as
Problem~\#243~\cite{erdosproblems}.

The additional hypothesis in Theorem~\ref{res:bounded} is the lower bound
on $E_n$. The growth and rationality assumptions already supply the other
state hypotheses after a finite shift. The nonnegative-error case, Badea's
criterion~\cite[Corollary~2.2, p.~316]{badea1993} in Koizumi's coordinates, is
integer descent \cite[Proposition~1(2), p.~14]{koizumi2025}; the theorem allows
arbitrary sign changes with bounded negative depth. The published
product-weighted criterion, stated as Koizumi's Corollary~4(1), assumes a
nonpositive upper limit \cite[pp.~14--15]{koizumi2025}, and
Corollary~\ref{res:originalbounded} allows any finite upper limit. Bado's
preprint, posted in September 2026, assumes two-sided bounded errors
\cite[Theorem~5.1, p.~4]{bado2026}; Theorem~\ref{res:bounded} needs only an
eventual lower bound.

\section{Proof of bounded-negative rigidity}\label{sec:bounded}
The arithmetic ingredients are proved immediately afterwards. Normalised
vanishing and integrality give $C_n\to\infty$ on a nonzero tail, while the
lower error bound gives $C_{n+1}\le C_n+B$.

\begin{proof}[Proof of Theorem~\ref{res:bounded}]
Suppose not. By Theorem~\ref{res:absorb}, after the centring threshold a
single zero would force every later error to be zero. Thus $E_n\ne0$
throughout a sufficiently late tail. If negative indices are infinite,
condition~(5) bounds their magnitudes. Proposition~\ref{res:gcdstab} makes
$G_n=\gcd(C_n,D_n)$ constant, say $g>0$, from some index onwards.
On this tail $u_n=C_n/g$ and $v_n=D_n/g$ form a reduced exact orbit.
Normalised vanishing and $|E_n|\ge1$ give $C_n\to\infty$, hence
$u_n\to\infty$. Also $u_{n+1}-u_n=-E_n/g\le B$.
Theorem~\ref{res:barrier} contradicts these bounded upward steps; when
$B=0$, descent alone already contradicts divergence.
If negative indices are finite, the remaining error is nonnegative and
Theorem~\ref{res:descent} makes it eventually zero, again a contradiction.
\end{proof}

\begin{corollary}\label{res:cor}
Under Theorem~\ref{res:bounded}, the multipliers satisfy
$a_{n+1}=a_n^2-a_n+1$ eventually.
\end{corollary}
\begin{proof}
Apply the two-zero implication in Corollary~\ref{res:eventual} below.
\end{proof}

\section{The arithmetic ingredients}\label{sec:state}
\subsection{Error, absorption and descent}\label{sec:defect}\label{sec:descent}
The identities below follow algebraically from Koizumi's state recurrence
\cite[Lemmas~2--4, pp.~10--12]{koizumi2025}. His Lemma~3 gives absorption,
and his Proposition~1(2) gives the nonnegative-error descent, credited there
to Badea \cite[Corollary~2.2, p.~316]{badea1993}.
The integer-remainder method for Cantor series is older
\cite[Theorem~2.1, pp.~85--86]{erdosstraus1974}. Han\v{c}l and Tijdeman
use a finite polynomial decomposition and the resulting shifted polynomial
sum to characterise rationality of polynomial Cantor series
\cite[Theorem~2.2 and its derivation, pp.~39--40]{hancltijdeman2008}.
Their polynomial hypotheses are essential to this comparison; we do not
invoke an unrestricted rearrangement of arrays of integer coefficients.
These are antecedents of the remainder method, not proofs of the
bounded-negative Chinese-remainder obstruction used here.

Write $\Delta_n=a_{n+1}-(a_n^2-a_n+1)$.

\begin{proposition}[error identities]\label{res:update}\label{res:defect}
For an exact integer state,
\[
 C_{n+1}=C_n-E_n,\qquad
 \Delta_n C_{n+1}=a_n^2E_n-E_{n+1}.
\]
\end{proposition}
\begin{proof}
Substitute $E_n=D_n-(a_n-1)C_n$ and use
$D_{n+1}=a_nD_n$, $C_{n+1}=a_nC_n-D_n$.
\end{proof}

\noindent Both identities are the
\lword{Erdos243/ReciprocalTailRigidity.lean}{57}{nextTailState_eq_sub_centered}{update law}
and the
\lword{Erdos243/ReciprocalTailRigidity.lean}{1775}{sylvesterDefect_mul_nextTailState}{defect identity}.

\begin{theorem}[absorption and descent]\label{res:absorb}\label{res:descent}
For a positive exact state with strict centring, $E_n=0$ implies
$E_{n+1}=0$. For any positive integer state with $C_{n+1}=C_n-E_n$,
eventual nonnegativity of $E_n$ implies its eventual vanishing.
\end{theorem}
\begin{proof}
When $E_n=0$, the second identity makes $E_{n+1}$ a multiple of
$C_{n+1}$; strict centring forces that multiple to be zero. In the second
assertion, $C_n$ is eventually a nonincreasing sequence of positive
integers, so it stabilises.
\end{proof}

\noindent Both assertions are
\lword{Erdos243/ReciprocalTailRigidity.lean}{2227}{centeredState_zero_absorbing}{absorption of a vanishing centred state}
and
\lword{Erdos243/ReciprocalTailRigidity.lean}{1843}{centeredState_eventually_zero}{stabilisation of the nonnegative error}.

\begin{corollary}[two zero errors]\label{res:step}\label{res:eventual}
If $E_n=E_{n+1}=0$ and $C_{n+1}\ne0$, then
$a_{n+1}=a_n^2-a_n+1$. Thus eventual zero error in a positive exact state
implies the eventual Sylvester recurrence.
\end{corollary}
\begin{proof}
The second error identity has nonzero factor $C_{n+1}$.
\end{proof}

\noindent The local and eventual forms are the
\lword{Erdos243/ReciprocalTailRigidity.lean}{1787}{sylvesterNext_eq_of_centered_zero}{local rigidity step}
and the
\lword{Erdos243/ReciprocalTailRigidity.lean}{1805}{sylvesterNext_eventually_of_centered_zero}{eventual Sylvester recurrence}.

Absorption holds after any eventual centring threshold. On a nonterminal
tail it removes all subsequent zeros. Then $|E_n|\ge1$, so normalised
vanishing gives $C_n\to\infty$.

\subsection{The first-crossing obstruction}\label{sec:barrier}
\begin{lemma}[consecutive multiples]\label{res:crt}
For pairwise coprime integers $m_0,\ldots,m_{B-1}\ge2$ and every lower
bound, there is a larger $t$ such that $m_i\mid t+i$ for each $i<B$.
\end{lemma}
\begin{proof}
Solve $t\equiv-i\pmod{m_i}$ by the Chinese remainder theorem, then add
multiples of $\prod_i m_i$.
\end{proof}

\noindent Lemma~\ref{res:crt} matches each modulus $m_i$ to its own multiple
$t+i$ inside a window of $B$ consecutive integers. The same forbidden block
and first-crossing contradiction, under a two-sided bound on the error, appear
in the proof of Bado's Theorem~5.1 \cite[pp.~4--5]{bado2026}. The lemma here
isolates bounded upward movement, and the gcd stabilisation below makes it
applicable with only a lower bound on the error.

\begin{theorem}[bounded-rise obstruction]\label{res:barrier}
Let $u:\N\to\N$ tend to infinity and let $B\ge1$ be an integer with
$u_{n+1}\le u_n+B$ for every $n$. There is no sequence of pairwise coprime
integers $m_i\ge2$ for which $\gcd(m_i,u_t)=1$ whenever $i<t$.
\end{theorem}
\begin{proof}
Choose $m_0,\ldots,m_{B-1}$ and use Lemma~\ref{res:crt} to find
$t>\max(u_0,\ldots,u_B)$ with $m_i\mid t+i$.  Since $u_n\to\infty$,
there is a first $n>B$ with $u_n\ge t$.  Minimality and the rise bound give
\[
 t\le u_n\le u_{n-1}+B<t+B.
\]
Hence $u_n=t+i$ for some $0\le i<B<n$, so $m_i\mid u_n$.
This contradicts $\gcd(m_i,u_n)=1$.  The argument uses first crossing,
not monotonicity: the numerator may fall before it reaches the block.
\end{proof}

\subsection{Reduction and stabilisation}
A reduced exact tail has positive $u_n$, integer $v_n\ge0$,
$\gcd(u_n,v_n)=1$, and
\[
 u_{n+1}=a_nu_n-v_n,\qquad v_{n+1}=a_nv_n.
\]
\begin{proposition}[persistent coprimality]\label{res:reduced}
In a reduced exact tail, $\gcd(a_n,v_n)=1$. Distinct multipliers are
pairwise coprime, and every earlier multiplier is coprime to every later
numerator.
\end{proposition}
\begin{proof}
A common prime divisor of $a_n,v_n$ would divide both $u_{n+1}$ and
$v_{n+1}$. Also, $a_i\mid v_t$ for $i<t$, so reducedness gives
$\gcd(a_i,u_t)=1$, and $\gcd(a_t,v_t)=1$ gives
$\gcd(a_i,a_t)=1$.
\end{proof}

\noindent All three conclusions are the
\lword{Erdos243/ReciprocalTailRigidity.lean}{989}{reducedStep_coprime_currentFactor}{step coprimality},
the
\lword{Erdos243/ReciprocalTailRigidity.lean}{1016}{reducedTail_pairwiseCoprime}{pairwise coprimality},
and the
\lword{Erdos243/ReciprocalTailRigidity.lean}{1030}{reducedTail_wholeModulusAvoidance}{whole-modulus avoidance}.

\begin{proposition}[gcd stabilisation]\label{res:gcdstab}
For a positive exact state, if negative errors have bounded magnitudes
along a cofinal set of indices, then $G_n=\gcd(C_n,D_n)$ eventually
stabilises. Division by its stable value gives a reduced exact tail.
\end{proposition}
\begin{proof}
Both updates preserve common divisors, so $G_n\mid G_{n+1}$. Moreover
$G_n\mid E_n$, hence $G_n\le -E_n$ at a negative index. The positive
divisibility chain is bounded along a cofinal set and therefore bounded
everywhere; it eventually stabilises. Division by the stable value
preserves both exact updates and leaves the states coprime.
\end{proof}

\noindent The stable gcd is
\lword{Erdos243/ReciprocalTailRigidity.lean}{1636}{tailGcd_eventuallyConstant_of_cofinally_boundedNegative}{gcd stabilisation}.

In the main proof, division by the stable gcd preserves divergence and a
bounded upward increment. Theorem~\ref{res:barrier} therefore applies to
the reduced numerator: this is the
\lword{Erdos243/ReciprocalTailRigidity.lean}{1066}{no_eventuallyBoundedRise_reducedTail}{eventual reduced exclusion}.
With normalised vanishing the same obstruction is the
\lword{Erdos243/ReciprocalTailRigidity.lean}{1754}{no_cofinallyBoundedNegative_of_normalizedVanishes}{exclusion from normalised vanishing}.
Sparse changes of the gcd without eventual
stabilisation are a different statement; the supporting record retains
that distinction and its finite-block consequences.

\section{Transfer to a reciprocal series}\label{sec:transfer}
The variables below are Koizumi's denominator-cleared rational tail
\cite[Lemma~4, pp.~11--12]{koizumi2025}, with the indexing and the common
clearing factor used here.
Let $x_n=\sum_{k\ge n}1/a_k$, suppose $x_1=p/q$ with $q>0$, and put
\[
 P_n=\prod_{j<n}a_j,\qquad D_n=qP_n,\qquad C_n=D_nx_n.
\]
Each $C_n$ is a positive integer: explicitly,
\[
 C_n=pP_n-q\sum_{k<n}P_n/a_k.
\]
Every quotient in the finite sum is an integer. Positivity follows from
$x_n>0$. The identity $x_n=1/a_n+x_{n+1}$ now gives both exact updates;
no approximation or rounding occurs in this construction.

For all large $n$, $a_{n+1}\ge a_n^2/2\ge2a_n$. Iterating the first
inequality after an index with $a_n>2$ gives
$a_n\ge\exp(c2^n)$ for some $c>0$. Summing the remaining doubled tail gives
\[
 x_n=\frac1{a_n}+\frac1{a_{n+1}}+O(a_{n+1}^{-2}).
\]
Consequently
\[
 \frac{C_{n+1}}{C_n}
 =\frac{a_nx_{n+1}}{x_n}
 =\frac{a_n^2}{a_{n+1}}+O(1/a_n)\longrightarrow1,
 \qquad \frac{E_n}{C_n}\longrightarrow0.
\]
Taking logarithms and averaging also gives $\log C_n=o(n)$; the same
argument gives $\log\max_{j\le n}C_j=o(n)$, since a fixed initial maximum
has negligible logarithm divided by $n$. Strict centring follows eventually. These are the ordinary
canonical-tail implications of Koizumi's work. After a finite shift the
sequence is the pseudo-greedy expansion of its own tail sum, with gap
sequence tending to zero \cite[Corollary~3, p.~9]{koizumi2025}. On that
tail the integers $c_n$, $d_n$ and $e_n$ of his Lemma~4
\cite[pp.~11--12]{koizumi2025} are $C_n$, $D_n$ and $E_n$ up to one common
positive factor, and his range $-c_n/2\le e_n<c_n/2$ gives $|E_n|<C_n$.

\begin{proof}[Proof of Corollary~\ref{res:originalbounded}]
Put $\gamma_n=a_n^2/a_{n+1}-1$. Since
\[
 \frac{P_{n+1}/a_{n+1}}{P_n/a_n}=1+\gamma_n\longrightarrow1,
 \qquad P_n/a_n=\exp(o(n)),
\]
the two-term estimate gives
\[
 \begin{aligned}
 E_n&=\frac{qP_n}{a_n}-\frac{qP_n(a_n-1)}{a_{n+1}}
 +O\!\left(\frac{qP_na_n}{a_{n+1}^2}\right),\\
 q\frac{P_n}{a_n}\gamma_n
 &=\frac{qP_na_n}{a_{n+1}}-\frac{qP_n}{a_n}.
 \end{aligned}
\]
Adding cancels the two terms of size $qP_n/a_n$ and leaves
\begin{equation}\label{eq:canonical-dictionary}
 E_n+q\frac{P_n}{a_n}\gamma_n
 =\frac{qP_n}{a_{n+1}}+
 O\!\left(\frac{qP_na_n}{a_{n+1}^2}\right)=o(1).
\end{equation}
Indeed, quadratic growth makes the two terms on the right respectively
$\exp(o(n))/a_n$ and $\exp(o(n))/a_n^2$, up to bounded factors.
The assumed upper bound therefore gives an eventual lower bound on the
integer $E_n$. After deleting a finite prefix, all hypotheses of
Theorem~\ref{res:bounded} hold. Apply Corollary~\ref{res:cor}.
\end{proof}

Identity~\eqref{eq:canonical-dictionary} is the product-weighted comparison
used in Koizumi's proof of his Corollary~4(1)
\cite[(11)--(12), p.~15]{koizumi2025}, written in the canonical variables.

The same estimate holds for any positive real scale $H_n\le qP_n$.
Put $W_n=H_nx_n$ and $Z_n=H_n-(a_n-1)W_n$; these quantities need not be
integers unless $H_n$ actually clears the denominator of $x_n$. Then
\begin{equation}\label{eq:general-clearance-dictionary}
 Z_n+\frac{H_n}{a_n}\gamma_n=o(1).
\end{equation}
For $H_n=L_n$ below, denominator clearing does hold, and the resulting
integral variables satisfy their own LCM update rather than the raw product
update. This supplies the LCM-prefactor consequence.

\section{A scalar finite-mass criterion}\label{sec:mass}
\begin{theorem}[finite negative relative mass]\label{res:massscalar}\label{res:mass}
Let $C_n$ be positive integers and $E_n$ integers satisfying
$C_{n+1}=C_n-E_n$. If
\[
 \sum_n\frac{(-E_n)_+}{C_n}<\infty,
\]
then $E_n=0$ eventually. Neither denominator dynamics nor normalised
vanishing is required. Positivity and integrality of $C_n$ are essential.
\end{theorem}
\begin{proof}
Set $\delta_n=(-E_n)_+/C_n$. Then
\[
 C_N\le C_0\prod_{n<N}(1+\delta_n)
 \le C_0\exp\!\left(\sum_n\delta_n\right).
\]
Choose an integer upper bound $K$ for $C_n$. Each strict rise contributes
at least $1/K$ to $\sum\delta_n$, so there are only finitely many rises.
The remaining positive integer sequence is nonincreasing and stabilises,
and the update then gives $E_n=0$.
\end{proof}

For comparison, the real sequence $C_n=1+1/(n+1)$ has
$E_n=C_n-C_{n+1}>0$ and zero negative mass, but never stabilises.
Thus the scalar conclusion is genuinely discrete.

\noindent The scalar statement, with no denominator dynamics or
normalised-vanishing assumption, is formalised as
\lword{Erdos243/SparseResetRecovery.lean}{612}{eventually_zero_of_summable_negativeRelativeMass_scalar}{vanishing from finite scalar negative mass}.
Its exact-orbit consequence is
\lword{Erdos243/SparseResetRecovery.lean}{690}{sylvesterNext_eventually_of_summable_negativeRelativeMass_scalar}{the Sylvester recurrence from finite scalar negative mass}.
The shorter proof above is included because this implication is an elementary
product bound followed by descent of positive integers.

On an exact reciprocal-tail orbit the conclusion gives the Sylvester
recurrence. For the gap sequence of the pseudo-greedy expansion, the same
criterion, with the same product bound and integer descent, appears in the
Erd\H{o}s Problem a Day working report on Problem~\#243, dated 12 August 2026
\cite[``A global termination criterion'']{erdosproblemaday243}; the statement
above is for any positive integer sequence with $C_{n+1}=C_n-E_n$.
The scalar criterion and bounded negative error are distinct
quantities to estimate; on the common exact-orbit class with normalised
vanishing they are each equivalent to the same endpoint.

\section{Weighted first crossings of LCM records}\label{sec:lcmrecords}

Only steps setting an LCM numerator record need contribute. The first-crossing
argument permits subtraction of a fixed baseline and any nonincreasing
nonnegative weight with divergent integral. We give the ordinary proof below.
The separate release also contains the full canonical weighted statement,
\href{https://github.com/wcook04/plectis-erdos-lean/blob/52f29ad173b04e3bac941b3663f2b9aebe5de0bb/ErdosProblems/Erdos243/PaperCompleteR8/CanonicalWeightedRecords.lean\#L268}{the real-weight criterion}; it is outside the checked main-repository build documented here.
The LCM coordinates also appear in Bado's September 2026 preprint:
with his denominator parameter equal to $q$, his $M_{n-1}$, $\Delta_{n-1}$,
$K_{n-1}$, $u_n$ and $g_n$ are the $L_n$, $M_n$, $U_n$, $V_n$ and $\rho_n$
below, and his (19) is $\rho_nU_{n+1}=U_n-V_n$
\cite[Proposition~7.1 and (16)--(20), pp.~6--7]{bado2026}.
This overlap concerns both the coordinates and their update, not merely
terminology.

Set
\[
 L_n=\operatorname{lcm}(q,a_1,\ldots,a_{n-1}),\quad
 M_n=D_n/L_n,\quad U_n=C_n/M_n,\quad V_n=E_n/M_n.
\]
The rational tail has denominator dividing $L_n$, so $U_n$ and $V_n$ are
integers. With $\rho_n=\gcd(L_n,a_n)$, the exact updates are
\[
 M_{n+1}=M_n\rho_n,\qquad
 \rho_nU_{n+1}=U_n-V_n,\qquad V_n=L_n-(a_n-1)U_n.
\]
Write $R_n=\max_{j\le n}U_j$ and
$\mathcal R=\{n:U_{n+1}>R_n\}$.  Strict centring already gives
$U_{n+1}<U_n$ when $\rho_n\ge2$, so every sufficiently late strict rise has
$\rho_n=1$.  The stronger eventual bound $-U_n\le2V_n$, supplied by
$V_n/U_n=E_n/C_n\to0$, gives the quantitative estimate
$U_{n+1}\le3U_n/4$ when $\rho_n\ge2$.
At a sufficiently late record step, where $\rho_n=1$, the actual jump is
$d_n=U_{n+1}-U_n=-V_n>0$. This identity is not asserted at a contracting
step with $\rho_n\ge2$.

\begin{theorem}[weighted record excess]\label{res:weightedrecord}
Assume the growth and rationality hypotheses of Problem~\ref{res:problem}.
Let $f:[1,\infty)\to[0,\infty)$ be finite and nonincreasing, with
$\int_1^\infty f(t)\,dt=\infty$. Then the sequence is eventually Sylvester
if and only if, for some integer $B\ge0$,
\[
 \sum_{n\in\mathcal R}(-V_n-B)_+f(U_n)<\infty.
\]
\end{theorem}

\begin{proof}
Suppose the sequence is not eventually Sylvester. Absorption and
integrality give $1/U_n\le |V_n|/U_n=|E_n|/C_n\to0$, so there are
infinitely many record steps. Here \emph{fresh} means $\rho_n=1$;
every sufficiently late record has this property. “Fresh” does not mean
that the digit itself is prime.
Their digits are pairwise coprime: an earlier digit divides the later
$L_n$, while the digit at a fresh record is coprime to $L_n$.
For any fixed $B\ge1$, choose $B$ such digits
$m_0,\ldots,m_{B-1}>B$ and take $T$ after their indices and after the
freshness threshold. Their size follows from $a_n\to\infty$.

Put $P=\prod_i m_i$ and choose $x$ by the Chinese remainder theorem with
$m_i\mid x+i$. Consider all translates $\tau=x+B+kP>R_T$, $k\in\Z$;
the first is at most $R_T+P$. A first crossing
$U_n\le R_n<\tau\le U_n+d_n$ is a record step. If $d_n\le B$, then
$U_n\in[\tau-B,\tau)$, so some $m_i$ divides $U_n$. It also divides
$L_n$, hence divides $d_n=(a_n-1)U_n-L_n$, contradicting
$0<d_n\le B<m_i$.

If the step first crosses $h\ge1$ such heights, their spacing gives
$(h-1)P<d_n$. With $r=d_n-B\ge1$ and $P\ge B+1$, we have
$d_n=B+r\le Pr$, whence $h\le r$. Monotonicity of $f$ now gives
\[
 \sum_{\substack{\tau\text{ first crossed}\\\text{at step }n}}f(\tau)
 \le(d_n-B)f(U_n).
\]
Each selected height above $R_T$ has exactly one first crossing. The first
selected height is at most $R_T+P$; each later height is $P$ further on.
For $R_N\ge R_T+P$, comparison of each interval with its left endpoint gives
\begin{equation}\label{eq:weightedcrossing}
 \sum_{\substack{T\le n<N\\n\in\mathcal R}}(-V_n-B)_+f(U_n)
 \ge\frac1P\int_{R_T+P}^{R_N} f(t)\,dt.
\end{equation}
Since $R_n\to\infty$, the right side diverges for every $B\ge1$;
$B=0$ follows by domination. Conversely a Sylvester tail telescopes to
$x_n=1/(a_n-1)$, so $V_n=0$ eventually.
\end{proof}

For example, $f(t)=1/[t\log(et)]$ gives the sufficient condition
\[
 \sum_{n\in\mathcal R}\frac{(-V_n-B)_+}{U_n\log(eU_n)}<\infty.
\]
Its finite lower bound in~\eqref{eq:weightedcrossing} is
$P^{-1}\log\bigl(\log(eR_N)/\log(e(R_T+P))\bigr)$.
Further fixed iterated logarithmic factors are allowed whenever the integral
still diverges. These are specialisations of one crossing theorem.

The criterion also has an exact expression in the original growth defect.
Put $\gamma_n=a_n^2/a_{n+1}-1$ and $\theta_n=E_n/C_n$. The defect identity
gives
\[
 \gamma_n+\theta_n=
 \frac{(1-\theta_n)(a_n-1+\theta_{n+1})}{a_{n+1}},
 \qquad 0<\gamma_n+\theta_n<3/a_n
\]
eventually. Thus the two nonnegative summands
$U_nf(U_n)(\gamma_n-B/U_n)_+$ and $(-V_n-B)_+f(U_n)$ differ by at most
$3U_nf(U_n)/a_n$. This is summable because
$U_n\le C_n=\exp(o(n))$, $f(U_n)\le f(1)$, and
$a_n\ge\exp(c2^n)$ eventually for some $c>0$.
Consequently Theorem~\ref{res:weightedrecord} is equivalent to finiteness of
\begin{equation}\label{eq:weightedgrowth}
 \sum_{n\in\mathcal R}U_nf(U_n)
 \left(\frac{a_n^2}{a_{n+1}}-1-\frac B{U_n}\right)_+
\end{equation}
for some $B$. The original hypotheses do not currently supply this
finiteness. In particular, termwise convergence to zero is insufficient.
Also, $d_n$ is the actual jump, including any recovery from a drawdown; it
must not be replaced by $R_{n+1}-R_n$ in the crossing proof.

\begin{corollary}[LCM-weighted bounded defect]\label{res:lcmbounded}
Assume the hypotheses of Problem~\ref{res:problem}. Write
$A_n=\operatorname{lcm}(a_1,\ldots,a_{n-1})$ with $A_1=1$ and
\[
 Z_n=\frac{A_n}{a_n}\Bigl(\frac{a_n^2}{a_{n+1}}-1\Bigr).
\]
If $\limsup Z_n<\infty$, then the sequence is eventually Sylvester.
\end{corollary}

\begin{proof}
Put $t_n=L_n/A_n=q/\gcd(q,A_n)$, so $1\le t_n\le q$.
Equation~\eqref{eq:general-clearance-dictionary}, with $H_n=L_n$, gives
$V_n+t_nZ_n=o(1)$. A finite upper bound on $Z_n$ therefore bounds the
negative part of $V_n$. The record series vanishes after a finite prefix
for a sufficiently large baseline, so Theorem~\ref{res:weightedrecord}
applies.
\end{proof}

Erd\H{o}s and Straus assume a nonpositive upper limit in the corresponding
next-index LCM expression \cite[Theorem~3, p.~132]{erdosstraus1964}, and
Tijdeman and Yuan give a criterion of the same kind for positive numerators
\cite{tijdemanyuan2002}. Here any finite upper
bound suffices under the quadratic-limit assumption. Since $A_n\mid P_n$,
this also implies Corollary~\ref{res:originalbounded}: multiplication by
a factor in $(0,1]$ preserves an upper bound, including at negative values.

Lean checks Corollary~\ref{res:lcmbounded} as
\href{https://github.com/wcook04/plectis-erdos/blob/3d6d938d696fed0fb71dd55115a18a73738ff223/lean/ErdosProblems/Erdos243/PaperCompleteR7/LcmDefect.lean#L49}{the original-coordinate LCM bounded defect}.
That declaration takes the rational reciprocal sum and the quadratic growth
limit as hypotheses, with the finite upper limsup recorded as an eventual
upper bound on $Z_n$; it assumes no divisibility of the denominator by the
prefix LCM.

Duverney proved a conditional signed form of the problem itself: for signs
$\epsilon_n\in\{-1,1\}$, if $\sum_{n\ge0}(a_{n+1}/a_n^{2}-1)$ converges,
then the reciprocal sum with numerators $\epsilon_n$ is rational if and
only if
\[
 a_{n+1}=a_n^2-(\epsilon_{n+1}/\epsilon_n)a_n
             +\epsilon_{n+2}/\epsilon_{n+1}
\]
for all large $n$~\cite[Corollary~3.2, p.~287]{duverney2001}. Absolute
convergence is a stronger sufficient specialisation, not the hypothesis
printed in that corollary. The all-positive specialisation is the form
relevant here.

For a quantitative comparison, Duverney, Kurosawa and Shiokawa
\cite[Theorem~1, author-version p.~2]{duverneykurosawashiokawa2020}
compute irrationality exponents for signed reciprocal series under eventual
$x_{n+1}\ge x_n^2$ and a denominator-growth condition on successive rational
ratios. The growth-side inequality excludes a Sylvester tail; that theorem is
not an unrestricted near-quadratic rigidity result.

The divisor constraint also survives integral numerator coefficients:
$d=(a_n-1)U_n-b_nL_n$ is divisible by every divisor of both $U_n,L_n$.
An ordinary coefficient-uniform variant gives bounded LCM height
from a bounded negative part; normalised vanishing is used afterwards for
stationarity. For positive numerator coefficients the classical comparators
are Badea's Corollary~2.2 \cite[p.~316]{badea1993} and the criterion of
Tijdeman and Yuan \cite{tijdemanyuan2002}. Its proof and exact bounded-height counterexamples are in the
\href{https://github.com/wcook04/plectis-erdos/blob/eccd8afc6db3c02a2265d0601b727d6c5e4467d5/lean/ErdosProblems/Erdos243/CoefficientUniformBoundedHeight.md}
{coefficient proof supplement}.

\section{Further consequences}\label{sec:secondaryrate}
\begin{corollary}[an inclusive one-sided $1/n$ bound]\label{res:inclusiveone}
Under the hypotheses of Corollary~\ref{res:originalbounded}, suppose that
for some $K\ge0$ and $\varepsilon>0$,
\[
 \gamma_n:=a_n^2/a_{n+1}-1\le \frac1n+\frac{K}{n^{1+\varepsilon}}
 \quad\hbox{for all large }n.
\]
Then the sequence is eventually Sylvester. In particular, the conclusion
holds under the pointwise eventual bound $\gamma_n\le1/n$.
\end{corollary}
\begin{proof}
Set $t_n=P_n/a_n$, so $t_{n+1}/t_n=1+\gamma_n$. For all large $n$,
\[
 1+\gamma_n\le(1+1/n)(1+K/n^{1+\varepsilon}).
\]
The first product telescopes and the second is bounded, giving $t_n=O(n)$.
Since $(\gamma_n)_+=O(1/n)$, the product defect $t_n\gamma_n$ is bounded
above. Apply Corollary~\ref{res:originalbounded}.
\end{proof}
This is an ordinary corollary of the bounded-defect theorem. It does not
assert that the weaker condition $\limsup n(\gamma_n)_+\le1$ suffices;
a vanishing but nonsummable excess over $1/n$ is not covered by this proof.

\begin{theorem}[cubic-rate irrationality]\label{res:cubicrate}
A strictly increasing sequence of positive integers with
\[
 a_n^2/a_{n+1}=1+\frac3n+o(n^{-3})
\]
has irrational reciprocal sum.
\end{theorem}
\begin{proof}[Proof by the polynomial exclusion]
Under rationality, Section~\ref{sec:transfer} gives
$C_{n+1}/C_n=1+3/n+o(n^{-3})$. Integer finite differences then give
$C_n=An(n+1)(n+2)+B$ eventually, with $A\in\Q_{>0}$ and $B\in\Q$.
The fixed-cubic exclusion contradicts this profile. The complete ordinary
argument, including the finite-difference extraction and the cubic-field
step, is in Section~2 of the
\href{https://wcook04.github.io/plectis/papers/erdos243-reciprocal-tail-reasoning-surface.pdf}{companion reasoning paper}.
The field step uses Chebotarev; see
\cite[Section~3, author version]{stevenhagenlenstra1996} for that classical
input. This paragraph is a proof by a stated companion result, not a
standalone finite-congruence proof and not an assembled Lean proof.
\end{proof}

The same finite-difference extraction excludes every nonintegral
$\lambda>1$ under the rate
$a_n^2/a_{n+1}=1+\lambda/n+o(n^{-\lambda})$: eventual polynomial growth
would force its degree to equal $\lambda$. The extraction uses
$\Gamma(n+\lambda)/\Gamma(n)\sim n^\lambda$
\cite[5.11.12]{dlmf_gamma}; its full proof is in the companion paper.
These rates fall outside $1+o(1/n)$, the rate at which
Koizumi notes that the Erd\H{o}s--Straus criterion settles the problem \cite[Remark~3, p.~16]{koizumi2025}, and they make
Duverney's signed series $\sum_n(a_{n+1}/a_n^{2}-1)$ diverge.

For $x_1=p/q$ the integer $C_n$ of Section~\ref{sec:transfer} is $q$ times
the linear form $P_nx_1-P_n\sum_{k<n}1/a_k$, whose coefficients are
integers. Approximation criteria for irrationality use nonzero forms of this
kind that tend to zero. At the cubic rate these forms grow: since $a_nx_n\to1$, the identity
$(P_{n+1}/a_{n+1})/(P_n/a_n)=1+\gamma_n$ of Section~\ref{sec:transfer} makes
$C_n/q=P_nx_n$ a positive multiple of $n(n+1)(n+2)$ up to a factor tending to one.
The proof of Theorem~\ref{res:cubicrate} uses the rate at precision
$o(n^{-3})$: for a rational sum it forces the exact profile
$C_n=An(n+1)(n+2)+B$, which is then excluded.

There is also an ordinary quantitative extension, not an end-to-end
Lean-checked theorem. Under the hypotheses of
Theorem~\ref{res:bounded} other than (5), put
$\operatorname{LL}(x)=\log_2\log_2\max(4,x)$. Each fixed
$\delta\in(0,1)$ gives
\[
 (-E_n)_+\le(1-\delta)\operatorname{LL}(C_n)
 \quad\text{eventually}\quad\Longrightarrow\quad E_n=0\text{ eventually}.
\]
Here $D_0>0$ follows after a shift from $a_n>1$, $C_n>0$ and normalised
vanishing. The LCM version needs the corresponding bound only at late record
steps.
The proofs use a CRT block in $[P,2P)$ and the canonical multiplier scale;
they are retained in
\href{https://github.com/wcook04/plectis-erdos/blob/eccd8afc6db3c02a2265d0601b727d6c5e4467d5/lean/ErdosProblems/Erdos243/SlowNegativePartRigidity.md}{the slow-negative proof}
and \href{https://github.com/wcook04/plectis-erdos/blob/eccd8afc6db3c02a2265d0601b727d6c5e4467d5/lean/ErdosProblems/Erdos243/LcmRecordExcess.md\#5-record-only-subcritical-log-log-bound}
{the record-only extension}. These are ordinary results with stated rate
hypotheses; neither establishes the general record budget.

Rapid decay alone should not be confused with the rigidity of this particular
recurrence. In the different setting of Erd\H{o}s Problem~\#270, Crmari\'{c}
and Kova\v{c} \cite[Theorems~1--2]{crmarickovac2025} show that the sums
$\sum_n\bigl(\prod_{j=1}^{f(n)}(n+j)\bigr)^{-1}$, with integer
$f(n)\to\infty$, range over all positive reals. With $f$ nondecreasing the
value set instead has measure zero, which does not establish irrationality
of every such sum. Their flexible-series construction is contextual here,
not an antecedent of the tail-state rigidity proof.

\section{The remaining arithmetic estimate}\label{sec:open}
\begin{proposition}[necessary profile]\label{res:frontier}
The canonical state of a sequence satisfying Problem~\ref{res:problem}'s
hypotheses but not its conclusion has $E_n\ne0$ eventually,
$|E_n|/C_n\to0$, unbounded negative magnitudes along negative indices,
and
\[
 \sum_n\frac{(-E_n)_+}{C_n}=\infty.
\]
\end{proposition}
\begin{proof}
Absorption excludes late zeros, descent excludes an eventually
nonnegative error, and Theorems~\ref{res:bounded} and~\ref{res:massscalar}
exclude the two finiteness conditions.
\end{proof}

\begin{problem}[global overlap-height growth]\label{res:lcmheight}
For every nonterminal canonical orbit satisfying Problem~\ref{res:problem},
must
\[
 \limsup_{n\to\infty}\frac{\log M_n}{n}>0?
\]
Equivalently, must there be a $K\ge1$ for which
\[
 2^n\le M_n^K
\]
at infinitely many indices?
\end{problem}
\noindent The necessary profile is proved above and is formalised as
\href{https://github.com/wcook04/plectis-erdos/blob/3d6d938d696fed0fb71dd55115a18a73738ff223/lean/ErdosProblems/Erdos243/PaperCompleteR7/Frontier.lean\#L151}{the canonical frontier theorem}. The formal statement assumes the rational reciprocal sum, the
quadratic growth limit and failure of eventual Sylvester behaviour; the
last hypothesis is essential to all of its nontermination conclusions.

For $B\ge0$ put
\[
 F_B(X)=\sum_{\substack{n\in\mathcal R\\U_n\le X}}(-V_n-B)_+.
\]
The remaining equivalent estimate is
\begin{equation}\label{eq:remaining-record-budget}
 \exists B\in\N:\qquad\liminf_{X\to\infty}\frac{F_B(X)}X=0.
\end{equation}
On a nonterminal orbit first crossings instead give
$F_B(X)\ge X/P_B-O_B(1)$, with an orbit-dependent CRT modulus $P_B$.
Charges are grouped by their source height $U_n$, not by their record height
$R_n$ or by their time index. Every fixed height has only finitely many late
sources because $U_n\to\infty$.
The missing step is to derive~\eqref{eq:remaining-record-budget} from
the canonical dynamics. Its equivalence with the existence of an admissible
weight is the following elementary lemma, applied to the record charges
$u_j=U_n$ and $w_j=(-V_n-B)_+$ over $n\in\mathcal R$ at a fixed $B$.

\begin{lemma}[weights and linear density]\label{res:weights}
Let $u_j$ be positive integers, let $w_j\ge0$, and put
$F(X)=\sum_{u_j\le X}w_j$, with the sum allowed a priori to be $+\infty$. Then $\liminf_{X\to\infty}F(X)/X=0$ if and only if
there is a finite nonincreasing $f:[1,\infty)\to[0,\infty)$ with
$\int_1^\infty f(t)\,dt=\infty$ and $\sum_jw_jf(u_j)<\infty$.
\end{lemma}

\begin{proof}
Suppose the lower limit vanishes. Since $F$ is constant on each $[m,m+1)$,
there are integers $X_1<X_2<\cdots$ with $X_k\ge2^{k}$ and
$F(X_k)\le2^{-k}X_k$. Put $f=\sum_{k\ge1}X_k^{-1}\mathbf 1_{[1,X_k]}$, which
is nonincreasing and bounded by $\sum_k2^{-k}$. Then
$\int_1^\infty f=\sum_k(1-1/X_k)=\infty$, while interchanging the nonnegative
sums gives $\sum_jw_jf(u_j)=\sum_kF(X_k)/X_k\le1$.

Conversely, let $f$ be such a weight and suppose $F(m)\ge cm$ for some $c>0$
and all integers $m\ge M$. A nonincreasing nonnegative function with divergent
integral is positive everywhere, so $F(X)\le f(X)^{-1}\sum_jw_jf(u_j)$ is
finite. Partial summation gives, for $N>M$,
\[
 \sum_{u_j\le N}w_jf(u_j)
 =F(N)f(N)+\sum_{m=1}^{N-1}F(m)\bigl(f(m)-f(m+1)\bigr)
 \ge c\Bigl(Mf(M)+\sum_{m=M+1}^{N}f(m)\Bigr),
\]
every term being nonnegative. The right side diverges as $N\to\infty$, because
$\sum_{m\ge1}f(m)\ge\int_1^\infty f(t)\,dt=\infty$.
\end{proof}

\noindent The lemma converts the weighted criterion into the displayed density
condition; supplying~\eqref{eq:remaining-record-budget} for a canonical orbit
remains open.

The obstruction is the exact unit-numerator feedback at the record sources.
The scalar sequence $C_n=n^2+1$, $E_n=-(2n+1)$ has normalised vanishing,
subexponential height and finite square mass, but is not an exact
reciprocal-tail orbit: its numerator word $0,0,2\pmod5$ violates the
persistent-zero transport of the exact equations. Likewise, small rises
and avoidance of an arbitrary sparse coprime family do not control the
canonical timing and size of the available divisors. The supporting record
keeps those falsifying examples with the exact hypotheses each preserves.
In particular, compatibility modulo one fixed old denominator does not
control later prime arrivals. The companion paper proves that arbitrary
finite words of units satisfy the reduced recurrence modulo that old
denominator, while separately deriving a large odd-prime-power supply from
the canonical growth. Neither fact alone supplies the missing record budget.

\paragraph{Formal scope and supporting record.}
At revision \texttt{\commitshort}, the supplied build receipt covers the
bounded-negative state theorem, its original-coordinate product- and
LCM-defect consequences, the scalar finite-mass theorem and the canonical
frontier theorem. A proof of a state implication is not, by itself, a proof of
its analytic premises; the canonical transfer is supplied in the separate
\texttt{PaperCompleteR7} modules cited above. The complete weighted criterion
has a proof source in the separate release, outside this main-repository
build. The cubic argument has formalised extraction, normalisation and finite
transport components, but no assembled declaration of the irrationality
theorem is identified here. The slow-negative results likewise distinguish
the checked arithmetic lemmas from the ordinary growth argument.
The accompanying reasoning record retains those distinctions, the finite
certificates and the unsuccessful extensions. No finite computation is used
as a substitute for a universal argument. The supply, valuation and energy
proofs added to the companion paper in this revision are ordinary proofs;
no new Lean replay is asserted. For older formal work on integer remainders,
see the Isabelle/HOL development of Koutsoukou-Argyraki and Li
\cite{kouli2020}; its endpoint is not Problem~\#243.

The authorship and AI-use statement on the first page applies to this
manuscript. Code and mathematical source are available at the repository
checkpoint \texttt{\commitshort}. A declaration establishes only its own
statement with its stated hypotheses. The finite upper bound and the global
record estimate remain additional mathematical obligations, not consequences
of a successful build.

\appendix
\section{Guide to the formal sources}\label{app:index}

Each linked phrase opens its declaration at the stated immutable revision.
The default checkpoint is \texttt{\commitshort}; separate-release and older
research-note links are explicitly identified by their destinations.
\texttt{ReciprocalTailRigidity.lean} contains the state dynamics and arithmetic
exclusions; \texttt{SparseResetRecovery.lean} contains the scalar finite-mass
criterion; and \texttt{PaperCompleteR7} supplies the original-coordinate
product-defect, LCM-defect and canonical-frontier theorems. Thus the
identification with reciprocal tails is not merely an informal interpretation.
The periodic exclusion still assumes $e_n<a_n$. The state theorem still lists
strict centring and normalised vanishing, with the former following from the
latter by taking $K=1$. Neither fact removes the additional bounded-negative
hypothesis. Section~\ref{sec:transfer} gives the ordinary analytic proof so the
reader need not inspect source code to follow the argument.

% The exhaustive source manifest is retained in the authored TeX for automated
% validation, and kept outside the printed note.
\iffalse

\begin{tabular}{@{}ll@{}}
informal statement & linked Lean source\\
Sylvester successor & \lrefx{Erdos243/ReciprocalTailRigidity.lean}{37}{sylvesterNext}\\
denominator update & \lrefx{Erdos243/ReciprocalTailRigidity.lean}{41}{nextDenState}\\
tail update & \lrefx{Erdos243/ReciprocalTailRigidity.lean}{45}{nextTailState}\\
centred state & \lrefx{Erdos243/ReciprocalTailRigidity.lean}{49}{centeredState}\\
Sylvester defect & \lrefx{Erdos243/ReciprocalTailRigidity.lean}{53}{sylvesterDefect}\\
update law & \lrefx{Erdos243/ReciprocalTailRigidity.lean}{57}{nextTailState_eq_sub_centered}\\
denominator equivariance & \lrefx{Erdos243/ReciprocalTailRigidity.lean}{65}{nextDenState_scale}\\
tail equivariance & \lrefx{Erdos243/ReciprocalTailRigidity.lean}{70}{nextTailState_scale}\\
centred equivariance & \lrefx{Erdos243/ReciprocalTailRigidity.lean}{75}{centeredState_scale}\\
normalised exclusion & \lrefx{Erdos243/ReciprocalTailRigidity.lean}{100}{no_normalizedConstantNegative_orbit}\\
scale-primitive exclusion & \lrefx{Erdos243/ReciprocalTailRigidity.lean}{147}{no_scalePrimitiveConstantNegative_orbit}\\
constant exclusion at every scale & \lrefx{Erdos243/ReciprocalTailRigidity.lean}{286}{no_constantNegative_orbit}\\
eventual constant exclusion & \lrefx{Erdos243/ReciprocalTailRigidity.lean}{350}{no_eventuallyConstantNegative_orbit}\\
persistence of a multiplier & \lrefx{Erdos243/ReciprocalTailRigidity.lean}{382}{base_dvd_denState_later}\\
one-step lock & \lrefx{Erdos243/ReciprocalTailRigidity.lean}{402}{divisorLocksTailSucc}\\
persistence of a lock & \lrefx{Erdos243/ReciprocalTailRigidity.lean}{429}{divisorLock_persists}\\
iterated period & \lrefx{Erdos243/ReciprocalTailRigidity.lean}{461}{periodic_add_mul}\\
iterated drift & \lrefx{Erdos243/ReciprocalTailRigidity.lean}{476}{phaseDrift_add_mul}\\
repeated-divisor transport & \lrefx{Erdos243/ReciprocalTailRigidity.lean}{494}{repeatedDivisor_forces_periodicCommonScale}\\
phase-primitive exclusion & \lrefx{Erdos243/ReciprocalTailRigidity.lean}{548}{no_phasePrimitivePeriodicNegative_orbit}\\
periodic exclusion at every scale & \lrefx{Erdos243/ReciprocalTailRigidity.lean}{690}{no_periodicNegative_orbit}\\
eventual periodic exclusion & \lrefx{Erdos243/ReciprocalTailRigidity.lean}{802}{no_eventuallyPeriodicNegative_orbit}\\
shifted consecutive multiples & \lrefx{Erdos243/ReciprocalTailRigidity.lean}{839}{exists_shifted_consecutiveMultiples}\\
bounded-rise barrier & \lrefx{Erdos243/ReciprocalTailRigidity.lean}{903}{no_boundedRise_of_tailAvoidance}\\
step coprimality & \lrefx{Erdos243/ReciprocalTailRigidity.lean}{989}{reducedStep_coprime_currentFactor}\\
pairwise coprimality & \lrefx{Erdos243/ReciprocalTailRigidity.lean}{1016}{reducedTail_pairwiseCoprime}\\
whole-modulus avoidance & \lrefx{Erdos243/ReciprocalTailRigidity.lean}{1030}{reducedTail_wholeModulusAvoidance}\\
reduced bounded-rise exclusion & \lrefx{Erdos243/ReciprocalTailRigidity.lean}{1041}{no_boundedRise_reducedTail}\\
eventual reduced exclusion & \lrefx{Erdos243/ReciprocalTailRigidity.lean}{1066}{no_eventuallyBoundedRise_reducedTail}\\
gcd monotonicity & \lrefx{Erdos243/ReciprocalTailRigidity.lean}{1430}{tailGcd_dvd_succ}\\
gcd at a negative index & \lrefx{Erdos243/ReciprocalTailRigidity.lean}{1587}{tailGcd_eq_gcd_negativeMagnitude}\\
chain stabilisation & \lrefx{Erdos243/ReciprocalTailRigidity.lean}{1605}{dvdChain_eventuallyConstant_of_cofinally_bounded}\\
gcd stabilisation & \lrefx{Erdos243/ReciprocalTailRigidity.lean}{1636}{tailGcd_eventuallyConstant_of_cofinally_boundedNegative}\\
tail-divergence form & \lrefx{Erdos243/ReciprocalTailRigidity.lean}{1661}{no_cofinallyBoundedNegative_of_tailTendsto}\\
divergence from normalised vanishing & \lrefx{Erdos243/ReciprocalTailRigidity.lean}{1739}{tailState_tendsto_atTop_of_nonzero_normalizedVanishes}\\
exclusion from normalised vanishing & \lrefx{Erdos243/ReciprocalTailRigidity.lean}{1754}{no_cofinallyBoundedNegative_of_normalizedVanishes}\\
defect identity & \lrefx{Erdos243/ReciprocalTailRigidity.lean}{1775}{sylvesterDefect_mul_nextTailState}\\
local rigidity step & \lrefx{Erdos243/ReciprocalTailRigidity.lean}{1787}{sylvesterNext_eq_of_centered_zero}\\
eventual Sylvester recurrence & \lrefx{Erdos243/ReciprocalTailRigidity.lean}{1805}{sylvesterNext_eventually_of_centered_zero}\\
eventual constancy & \lrefx{Erdos243/ReciprocalTailRigidity.lean}{1825}{antitone_nat_eventually_constant}\\
stabilisation of the nonnegative error & \lrefx{Erdos243/ReciprocalTailRigidity.lean}{1843}{centeredState_eventually_zero}\\
tail realisation & \lrefx{Erdos243/ReciprocalTailRigidity.lean}{1869}{natTail_eq_nextTailState}\\
denominator realisation & \lrefx{Erdos243/ReciprocalTailRigidity.lean}{1881}{natDen_eq_nextDenState}\\
signed update law & \lrefx{Erdos243/ReciprocalTailRigidity.lean}{1978}{natTail_eq_sub_centeredState}\\
absorption of a vanishing centred state & \lrefx{Erdos243/ReciprocalTailRigidity.lean}{2227}{centeredState_zero_absorbing}\\
contrapositive along a tail & \lrefx{Erdos243/ReciprocalTailRigidity.lean}{2254}{eventually_nonzero_of_zero_absorbing}\\
bounded-negative-part rigidity & \lrefx{Erdos243/ReciprocalTailRigidity.lean}{2271}{boundedNegativePart_eventually_zero}\\
eventual bounded-negative-part rigidity & \lrefx{Erdos243/ReciprocalTailRigidity.lean}{2360}{eventuallyBoundedNegativePart_eventually_zero}\\
repair-entropy inequality & \lrefx{Erdos243/RepairEntropy.lean}{48}{repairedFamily_recovery_energy_divisionFree}\\
independent-repair bound & \lrefx{Erdos243/RepairEntropy.lean}{78}{repaired_card_bound_of_independent}\\
fixed-length reset disappearance & \lrefx{Erdos243/RepairEntropy.lean}{107}{eventually_recoveryPayment_eq_one_of_fixedLength}\\
forced numerator & \lrefx{Erdos243/FiniteHorizonResidue.lean}{22}{forcedNumerator}\\
polynomial congruence & \lrefx{Erdos243/FiniteHorizonResidue.lean}{26}{forcedNumerator_modEq}\\
exact-division cancellation & \lrefx{Erdos243/FiniteHorizonResidue.lean}{41}{quotient_modEq_of_modEq_mul}\\
horizon modulus & \lrefx{Erdos243/FiniteHorizonResidue.lean}{53}{horizonModulus}\\
ascending-factorial form & \lrefx{Erdos243/FiniteHorizonResidue.lean}{59}{horizonModulus_eq_ascFactorial}\\
factorial value at the initial index & \lrefx{Erdos243/FiniteHorizonResidue.lean}{69}{horizonModulus_zero_eq_factorial}\\
survival predicate & \lrefx{Erdos243/FiniteHorizonResidue.lean}{75}{ForcedSurvives}\\
quotient transport & \lrefx{Erdos243/FiniteHorizonResidue.lean}{85}{forcedQuotient_modEq}\\
shrinking-modulus transport & \lrefx{Erdos243/FiniteHorizonResidue.lean}{97}{forcedSurvives_iff_of_modEq}\\
factorial residue reduction & \lrefx{Erdos243/FiniteHorizonResidue.lean}{134}{forcedSurvives_iff_of_modEq_factorial}\\
\end{tabular}
\fi

\section{A factorial residue reduction for forced orbits}\label{app:residue}

A constant-negative error $E_n=-m$ with $C_0=c$ produces $C_n=c+nm$ and the
shape equation
\begin{equation}\label{eq:shape}
 D_n+m=(a_n-1)C_n.
\end{equation}
In the case $m=c=1$, where~\eqref{eq:shape} reads $D_n+1=(a_n-1)(n+1)$, each
multiplier is determined by its predecessor; we call such an orbit
\emph{forced}.  At index $n$ the numerator of the next multiplier is
\[
 \fnum(n,a)=(n+1)a^{2}-(n+2)a+(n+3),
\]
the
\lword{Erdos243/FiniteHorizonResidue.lean}{22}{forcedNumerator}{forced
numerator}, and the divisor is $n+2$.  The orbit survives a step when that
division is exact, giving a survival predicate, the
\lword{Erdos243/FiniteHorizonResidue.lean}{75}{ForcedSurvives}{survival
predicate}.
\begin{example}[the shape equation running until it fails]\label{ex:shape}
The forced orbit from $a=3$ reads this way: $\fnum(0,3)=6$ is divisible by $2$
and gives $a_1=3$, while $\fnum(1,3)=13$ is not divisible by $3$, so the orbit
stops there.
\end{example}
Deciding survival by iteration is expensive because the orbit
grows doubly exponentially, and it is unnecessary: survival over a finite
horizon depends on the initial value only through a factorial residue.
Computing a pseudo-greedy orbit through residues modulo a shrinking product
modulus is Koizumi's method \cite[Remark~2 and Algorithm~1,
pp.~13--14]{koizumi2025}; for the forced numerator that product is a
factorial.

\begin{theorem}[factorial residue reduction]\label{res:residue}
For all $h$ and all integers $a\equiv b \pmod{(h+1)!}$, the orbit from $a$
survives $h$ forced updates exactly when the orbit from $b$ does.
\end{theorem}

\begin{proof}
Let $M(0,i)=1$ and $M(h+1,i)=(i+2)M(h,i+1)$, an ascending factorial with
$M(h,0)=(h+1)!$.  The numerator is a polynomial with integer coefficients, so
congruences transfer; reducing the modulus gives divisibility by $i+2$ for one
exactly when for the other, and cancelling that common factor from both values
and modulus leaves the inductive hypothesis at $M(h,i+1)$.
\end{proof}

\noindent Formalised as the
\lword{Erdos243/FiniteHorizonResidue.lean}{134}{forcedSurvives_iff_of_modEq_factorial}{factorial
residue reduction}, over the
\lword{Erdos243/FiniteHorizonResidue.lean}{97}{forcedSurvives_iff_of_modEq}{shrinking-modulus
transport}, the
\lword{Erdos243/FiniteHorizonResidue.lean}{26}{forcedNumerator_modEq}{polynomial
congruence}, and the
\lword{Erdos243/FiniteHorizonResidue.lean}{41}{quotient_modEq_of_modEq_mul}{exact-division
cancellation}; the modulus identifications are the
\lword{Erdos243/FiniteHorizonResidue.lean}{59}{horizonModulus_eq_ascFactorial}{ascending-factorial
form} and the
\lword{Erdos243/FiniteHorizonResidue.lean}{69}{horizonModulus_zero_eq_factorial}{factorial
value at the initial index}.

At $h=1$ the modulus is $2!=2$, and surviving one update means
$2\mid a^{2}-2a+3$, which holds exactly for odd $a$: for instance
$\fnum(0,3)=6$ but $\fnum(0,4)=11$.  So one step of survival is decided by the
parity of $a$ alone, which is Theorem~\ref{res:residue} at its smallest
nontrivial horizon.

The inherited search record reports a maximum forced-prefix length of $17$
among initial states below $5000$. This finite search was not rerun for the
present revision and is not a proof of exclusion; the Lean-checked
\lword{Erdos243/ReciprocalTailRigidity.lean}{286}{no_constantNegative_orbit}{constant exclusion at every scale}
now excludes the constant-negative case outright, for every seed and at every
scale.  The reduction is retained because it is exact, and because the
shrinking-modulus technique transfers to any forced orbit whose step is a
polynomial division.

\begin{thebibliography}{99}
\small
\raggedright
\setlength{\itemsep}{1pt}
\bibitem{erdosstraus1964} P.~Erd\H{o}s and E.~G. Straus,
  \href{https://users.renyi.hu/~p_erdos/1964-19.pdf}{\emph{On the irrationality
  of certain Ahmes series}},
  J. Indian Math. Soc. (N.S.) \textbf{27} (1964), 129--133. MR~175848.
\bibitem{duverney2001} D.~Duverney,
  \href{https://www.ms.u-tokyo.ac.jp/journal/pdf/jms080206.pdf}
  {\emph{Irrationality of fast converging series of rational numbers}},
  J. Math. Sci. Univ. Tokyo \textbf{8} (2001), 275--316. MR~1837165.
\bibitem{badea1993} C.~Badea,
  \href{https://matwbn.icm.edu.pl/ksiazki/aa/aa63/aa6342.pdf}{\emph{A theorem
  on irrationality of infinite series and applications}}, Acta Arith.
  \textbf{63} (1993), no.~4, 313--323,
  doi:\href{https://doi.org/10.4064/aa-63-4-313-323}{10.4064/aa-63-4-313-323}.
\bibitem{tijdemanyuan2002} R.~Tijdeman and P.~Yuan,
  \emph{On the rationality of Cantor and Ahmes series},
  Indag. Math. (N.S.) \textbf{13} (2002), no.~3, 407--418,
  doi:\href{https://doi.org/10.1016/S0019-3577(02)80018-0}{10.1016/S0019-3577(02)80018-0}.
\bibitem{erdosgraham1980} P.~Erd\H{o}s and R.~L. Graham,
  \href{https://mathweb.ucsd.edu/~ronspubs/80_11_number_theory.pdf}{\emph{Old
  and New Problems and Results in Combinatorial Number Theory}},
  Monogr. Enseign. Math. 28, Geneva, 1980, p.~64.
\bibitem{erdos1988} P.~Erd\H{o}s,
  \href{https://users.renyi.hu/~p_erdos/1988-22.pdf}{\emph{On the irrationality
  of certain series: problems and results}}, in A.~Baker (ed.), \emph{New
  Advances in Transcendence Theory}, Cambridge UP, 1988, pp.~102--109,
  doi:\href{https://doi.org/10.1017/CBO9780511897184.009}{10.1017/CBO9780511897184.009}.
\bibitem{bado2026} I.~O. Bado,
  \emph{Prime-Support Rigidity and Primitive Pseudo-Greedy Dynamics:
  Partial Progress on Erd\H{o}s Problem~\#243}, preprint posted September 2026,
  doi:\href{https://doi.org/10.13140/RG.2.2.36612.08325}
  {10.13140/RG.2.2.36612.08325}.
\bibitem{erdosproblemaday243} P.~White with Claude (Anthropic),
  \href{https://erdosproblemaday.com/report/243}{\emph{Erd\H{o}s \#243:
  working report}}, Erd\H{o}s Problem a Day, page dated 12 August 2026.
  AI-assisted, unrefereed working report.
\bibitem{stevenhagenlenstra1996} P.~Stevenhagen and H.~W. Lenstra, Jr.,
  \href{https://pub.math.leidenuniv.nl/~lenstrahw/PUBLICATIONS/1994c/art.pdf}
  {\emph{Chebotar\"ev and his density theorem}}, Math. Intelligencer
  \textbf{18} (1996), no.~2, 26--37,
  doi:\href{https://doi.org/10.1007/BF03027290}{10.1007/BF03027290}.
  Page references are to the linked author version dated 23 March 1995.
\bibitem{koizumi2025} J.~Koizumi,
  \href{https://math.colgate.edu/~integers/aa28/aa28.pdf}
  {\emph{Irrationality of the reciprocal sum of doubly exponential sequences}},
  Integers \textbf{26} (2026), Paper No.~A28, 17 pp.,
  doi:\href{https://doi.org/10.5281/zenodo.18714404}{10.5281/zenodo.18714404}.
  Locators refer to this published version.
\bibitem{erdosproblems} T.~F. Bloom,
  \href{https://www.erdosproblems.com/243}{\emph{Erd\H{o}s Problem \#243}},
  supplied snapshot of 28 July 2026; not a live status verification.
\bibitem{erdosstraus1974} P.~Erd\H{o}s and E.~G. Straus,
  \href{https://msp.org/pjm/1974/55-1/pjm-v55-n1-p08-s.pdf}
  {\emph{On the irrationality of certain series}}, Pacific J. Math.
  \textbf{55} (1974), no.~1, 85--92.
\bibitem{hancltijdeman2008} J.~Han\v{c}l and R.~Tijdeman,
  \href{https://www.impan.pl/shop/en/publication/transaction/download/product/82186}
  {\emph{On the irrationality of polynomial Cantor series}},
  Acta Arith. \textbf{133} (2008), no.~1, 37--52,
  doi:\href{https://doi.org/10.4064/aa133-1-3}{10.4064/aa133-1-3}.
  Locators refer to the published version.
\bibitem{duverneykurosawashiokawa2020} D.~Duverney, T.~Kurosawa and I.~Shiokawa,
  \href{https://danielduverney.fr/documents/theorie-des-nombres/Tsukuba.pdf}
  {\emph{Irrationality exponents of certain fast converging series of rational
  numbers}}, Tsukuba J. Math. \textbf{44} (2020), no.~2, 235--250,
  doi:\href{https://doi.org/10.21099/tkbjm/20204402235}{10.21099/tkbjm/20204402235}.
  Theorem locators follow the linked 14-page author version.
\bibitem{crmarickovac2025} T.~Crmari\'{c} and V.~Kova\v{c},
  \href{https://arxiv.org/abs/2504.18712v1}
  {\emph{On the irrationality of certain super-polynomially decaying series}},
  Colloq. Math. \textbf{179} (2025), 55--68,
  doi:\href{https://doi.org/10.4064/cm9628-5-2025}{10.4064/cm9628-5-2025}.
  Theorem locators follow arXiv:2504.18712v1.
\bibitem{kouli2020} A.~Koutsoukou-Argyraki and W.~Li,
  \href{https://isa-afp.org/entries/Irrational_Series_Erdos_Straus.html}
  {\emph{Irrationality Criteria for Series by Erd\H{o}s and Straus}},
  Archive of Formal Proofs, 12 May 2020. An Isabelle/HOL formalisation;
  the archive entry identifies the results formalised.
\bibitem{dlmf_gamma} National Institute of Standards and Technology,
  \href{https://dlmf.nist.gov/5.11.E12}{\emph{Digital Library of Mathematical
  Functions}, \S5.11(iii), formula 5.11.12}, accessed 16 September 2026.
\end{thebibliography}

\companionpapercontext

\end{document}
